% ==============================================================
% metadata.tex — Article
% Metadata for a multi-author academic article. Authors are listed
% manually (AuthorA, AuthorB, AuthorC by default) rather than via a
% loop or Python wrapper -- add/remove blocks here and mirror the
% change in partials/cover.tex and definitions/headerfooter.tex.
% ==============================================================

% --------------------------------------------------------------
% Draft mode
\newcommand{\IsDraft}{false}

% --------------------------------------------------------------
% Content
\newcommand{\DocType}{Research Article}
\newcommand{\DocTitle}{[Article Title]}
\newcommand{\DocJournal}{[Journal Title]}
\newcommand{\DocAbstract}{\lipsum[1] \lipsum[2]}
\newcommand{\DocAcknowledgements}{\lipsum[1]}
\newcommand{\DocKeywords}{[Keyword One]; [Keyword Two]; [Keyword Three]}
\newcommand{\DocHighlightOne}{\lipsum[1][1-2]}
\newcommand{\DocHighlightTwo}{\lipsum[2][1-2]}
\newcommand{\DocHighlightThree}{\lipsum[3][1-2]}

% --------------------------------------------------------------
% Authorship
% Default: 3 authors. Each carries name, contact email, a
% corresponding-author flag ("true"/"false", read via \ifdefstring),
% and an affiliation number matching an entry in \DocAffiliations
% below. Two authors at the same institution should share a number
% (e.g. set \AuthorBAffiliationNum to {1} to match Author A) -- when
% you do that, also delete that author's now-duplicate entry from
% \DocAffiliations by hand.
% Each author also carries a \AuthorXCredit field: a manually written,
% comma-separated list of CRediT (Contributor Roles Taxonomy) roles --
% e.g. Conceptualization, Data Curation, Formal Analysis, Funding
% Acquisition, Investigation, Methodology, Project Administration,
% Resources, Software, Supervision, Validation, Visualization,
% Writing -- Original Draft, Writing -- Review \& Editing.
\newcommand{\AuthorAName}{[Author A First Name]}
\newcommand{\AuthorALastName}{[Author A Last Name]}
\newcommand{\AuthorAAffiliationNum}{1}
\newcommand{\AuthorAEmail}{authora@example.com}
\newcommand{\AuthorACorresponding}{true}
\newcommand{\AuthorACredit}{Conceptualization, Methodology, Writing -- Original Draft}

\newcommand{\AuthorBName}{[Author B First Name]}
\newcommand{\AuthorBLastName}{[Author B Last Name]}
\newcommand{\AuthorBAffiliationNum}{2}
\newcommand{\AuthorBEmail}{authorb@example.com}
\newcommand{\AuthorBCorresponding}{false}
\newcommand{\AuthorBCredit}{Data Curation, Formal Analysis, Investigation}

\newcommand{\AuthorCName}{[Author C First Name]}
\newcommand{\AuthorCLastName}{[Author C Last Name]}
\newcommand{\AuthorCAffiliationNum}{3}
\newcommand{\AuthorCEmail}{authorc@example.com}
\newcommand{\AuthorCCorresponding}{false}
\newcommand{\AuthorCCredit}{Software, Visualization, Writing -- Review \& Editing}

% --------------------------------------------------------------
% Affiliations
% One semicolon-separated list, numbered to match the
% \AuthorXAffiliationNum values above. Edit this single line by hand
% to add, remove, reorder, or merge entries.
\newcommand{\DocAffiliations}{1 [Author A Affiliation]; 2 [Author B Affiliation]; 3 [Author C Affiliation]}

% --------------------------------------------------------------
% Corresponding author
% Point this at whichever AuthorX has \AuthorXCorresponding set to
% {true} above -- kept as a separate manual field rather than derived.
\newcommand{\DocCorrespondingAuthorEmail}{\AuthorAEmail}

% --------------------------------------------------------------
% Identifiers
\newcommand{\DocYear}{\the\year}
\newcommand{\DocDOI}{https://doi.org/10.0000/xxxxxx}

% --------------------------------------------------------------
% Draft watermark
\ifdefstring{\IsDraft}{true}{
    \usepackage{draftwatermark}
}{}
